% Responsable: A
\section{Introducción}
\subsection{Contexto}

La reducción de la mortalidad materna y fetal, junto con tratamientos oportunos
durante el embarazo son factores prioritarios a nivel de sanidad pública global
lo cual se alinea con los ODS, como se menciona en ONU Mujeres ``\emph{Meta 3.1
de los ODS: Reducir la tasa mundial de mortalidad materna a menos de 70 por cada
100\,000 nacidos vivos\ldots Un poco más del 90\,\% de todas las muertes maternas
ocurren en países de ingresos y medios bajos en 2023}''. Por eso el principal
problema radica en situaciones de monitorización primaria en zonas rurales o
vulnerables de un país en desarrollo junto con la identificación temprana de los
síntomas y/o complicaciones para descubrir el nivel de riesgo en función del
seguimiento continuo de signos vitales como la frecuencia cardíaca, presión
arterial sistólica, temperatura corporal, entre otras. Sin embargo, evaluaciones
puntuales o un control inconsistente dificultan la detección a tiempo de
condiciones médicas que afectan tanto a la madre como al bebé y la Organización
Mundial de la Salud [OMS] (2025) indica que \emph{la mayoría de las muertes
maternas son prevenibles, ya que existen soluciones sanitarias bien conocidas
para prevenir o controlar las complicaciones}.

No obstante, se depende demasiado del presupuesto público y gubernamental para el
financiamiento de proyectos de tecnologías o nuevas infraestructuras en la salud
pública (Ramos Vázquez, 2019), ``\emph{La mayoría de muertes, lesiones y
problemas de salud se pueden evitar mejorando el acceso a servicios adecuados de
atención de la salud, incluidos los métodos de planificación familiar eficaces y
la atención obstétrica de emergencia, reconociendo el derecho de la mujer y del
hombre a la información y al acceso a métodos seguros, eficaces y
asequibles\ldots}''. Bajo esta perspectiva y considerando que las complicaciones
no detectadas a tiempo generan hospitalizaciones obstétricas, esto representa un
costo significativamente mayor para el sistema de salud y las familias.

Asimismo, existen grandes diferencias respecto a la atención médica prenatal
ofrecida o disponible en los distintos centros de atención rurales, alejados de
las ciudades donde hay una brecha sea tecnológica o económica que dificulta el
debido proceso de monitorización durante el embarazo. Además, como se mencionó
anteriormente, según la ONU el difícil acceso a recursos, bienes y servicios del
sistema de saneamiento básico en ciertas comunidades vulnerables puede
incrementar la susceptibilidad a infecciones durante la gestación y en
consecuencia enfermedades (OMS, 2025).

Dicho esto y en síntesis de lo anterior, el cumplimiento estricto de las
regulaciones sobre privacidad y confidencialidad de la información médica de los
pacientes es fundamental al igual que las exigencias reglamentarias para la
validación y certificación de modelos de Inteligencia Artificial y/o Estadística
aplicados a la salud humana (Proyecto de Ley 216, 2026), dado a los tipos de
datos a tratar, en este caso, son datos recopilados de diferentes hospitales,
centros médicos de salud materna y clínicas comunitarias de las zonas rurales de
Bangladesh mediante un sistema de monitoreo de riesgos basado en IoT, con 1014
registros de pacientes con variaciones simultáneas en distintos indicadores
fisiológicos, lo que requiere procesos de identificación de factores de riesgo
para la adecuada intervención médica.

\subsection{Planteamiento del problema}

No se ha profundizado lo suficiente en la evaluación de la suficiencia de los
datos clínicos de rutina para la detección oportuna de riesgos durante el
embarazo en entornos con recursos limitados. Aunque la literatura indica que la
información clínica de los controles prenatales posee un valor predictivo mayor
respecto a los datos demográficos, la necesidad de determinar qué tan viable y
preciso resulta implementar estos modelos predictivos cuando se trabaja bajo la
restricción de muestras pequeñas y la escasez de personal u equipo médico aún se
evidencia.

\subsection{Pregunta de investigación}

¿En qué medida las variables fisiológicas registradas (presión arterial,
glucosa, temperatura corporal, frecuencia cardíaca y edad) permiten clasificar
con precisión el nivel de riesgo (RiskLevel) de una gestante, y cuáles de estos
factores tienen mayor peso predictivo?

\subsection{Justificación}

El proyecto nos ofrecerá un sistema predictivo de triaje prenatal adaptado a
entornos de atención primaria en comunidades rurales de escasos recursos. Con
una de las bases, las sugerencias de optimización de hiperparámetros y selección
de modelos (Nogholi et al., 2026), se espera construir un modelo
computacionalmente eficaz y de alto rendimiento que priorice la sensibilidad o
\emph{recall}, garantizando la detección oportuna de gestantes en estado crítico
antes del parto. Este resultado permitirá aclarar la toma de decisiones clínicas
y optimizar la asignación de recursos públicos en zonas rurales, orientando el
transporte médico como ambulancias y la remisión a centros hospitalarios
únicamente a los casos que representan un verdadero alto riesgo y no una falsa
alarma (Nogholi et al., 2026).

Por otro lado, los principales riesgos del proyecto serían el sobreajuste o la
pérdida de precisión al transferir el modelo a otras regiones desatendidas y/o
contextos (Nogholi et al., 2026), la presencia de ruido o fallas en la captura
de características fisiológicas de las madres gestantes mediante sensores IoT, y
al ser modelos opacos (dan el resultado pero no una explicación médica al
respecto), se genera una falta de interpretación clínica y desconfianza del
personal médico. Al no integrar mecanismos de Inteligencia Artificial Explicable
(XAI), existe el peligro de que la herramienta sea ignorada en la práctica real
o rechazada (Nogholi et al., 2026).

\subsection{Objetivo general}

Desarrollar un modelo de clasificación que, a partir de las variables
fisiológicas registradas (presión arterial sistólica y diastólica, glucosa,
temperatura corporal, frecuencia cardíaca y edad), permita establecer con
precisión el nivel de riesgo (RiskLevel) de una gestante, identificando los
factores con mayor peso predictivo para facilitar la intervención médica
oportuna.

\subsection{Objetivos específicos}

\begin{enumerate}
  \item Describir el comportamiento y distribución de las variables fisiológicas
        según el nivel de riesgo.
  \item Identificar los factores fisiológicos más asociados con un riesgo alto.
  \item Construir y comparar modelos que permitan clasificar el nivel de riesgo
        con mayor precisión.
\end{enumerate}

% TODO: contexto, planteamiento del problema, pregunta de investigación,
% justificación, objetivo general y específicos, estructura del documento.
